% Intro — full prose v1 (written last-but-one; abstract polished after).

Two appealing beliefs circulate about the reasoning traces of large
language models. The first is that traces contain \emph{breakthrough
moments} --- points where accumulated reasoning suddenly unlocks a
solution, popularized as ``aha moments'' in reinforcement-trained
reasoners \citep{deepseekr1}. The second is that a run's fate is legible
early: probes on hidden states reportedly predict correctness
with AUROC $0.79$--$0.95$, sometimes from the first four tokens
\citep{zhang2025knowright, yuan2026diagnostic, david2025temporal,
sun2026trajectories}, and
serving systems already allocate compute on such signals
\citep{fu2024certaindex, fu2025deepconf}. Both beliefs, if true, would
matter: the first for what test-time compute buys, the second for
adaptive inference and interpretability.

Both beliefs rest on measurements with a missing control, and the two
missing controls mirror each other: a claim about a trajectory requires a
counterfactual control at the level of the claim --- what this
computation buys over fresh computation, what this attempt reveals beyond
its problem. Breakthrough claims rest on
truncate-and-resample probes: cut the trace at position $t$, sample
continuations under a budget, and call the first high-solve-rate anchor a
breakthrough. But a crossing of that kind conflates two events ---
\emph{the prefix accumulated value} and \emph{the remaining work began to
fit the budget}. Distinguishing them requires a from-scratch \emph{restart}
control at matched total generated-token budget; we found no prior work
combining per-anchor continuation estimates with a per-problem restart
curve at matched budget (Section~\ref{sec:related}).
Prediction claims rest on probes evaluated across problems, where a probe
can score well by reading \emph{which problem it is} --- its difficulty
--- rather than how the attempt is going. Distinguishing those requires a
question-only baseline and a within-problem evaluation, which the
published positives do not report.

This paper supplies both controls and reports what they change. We build a
restart-controlled, budget-indexed, censoring-aware truncation probe
(Section~\ref{sec:method}): per-anchor continuation solve rates
$\phat(t)$ are compared against restart curves $R(C)$ at matched total
generated-token budget, splitting the naive breakthrough time into a budget-fit time
\Tf{} and a prefix-value time \Tv, with attempt-level noise hardened by
replication and pooling. We apply it to 178 problem--model cells (89 MATH
problems $\times$ two small open models, 16K-token instrumented traces),
and we complement it with two analyses of public reasoning-model corpora
at larger scale.

The instrument returns a consistent picture:

\begin{itemize}
\item \textbf{Breakthroughs are mostly budget artifacts.} Exactly 1 of 178
cells survives as prefix-limited (three cross the advantage margin; an
in-grid restart already solves two); 98 cells collected after the
rules were frozen contributed zero new cases
(Section~\ref{sec:results-regimes}).
\item \textbf{Restart dose--response separates failure modes.} One model's
unsolved problems dissolve as budget grows ($0\%\to79\%$,
compute-starved); the other's do not budge (capability-limited within
the measured restart-budget range).
\item \textbf{Accumulated reasoning is predominantly compute
compression.} At matched total token budget the model's own prefix beats
restarting wherever the comparison is exactly matched (9/9, with 2 wins
and 2 ties on boundary proxies), yet a larger restart budget reaches the
same threshold in 11/13 and the prefix's own rate in 9/13: within the
measured budget range, long reasoning mostly buys the same successes
cheaper rather than reaching otherwise unreachable ones.
\item \textbf{Solvability is not binary.} A third of ambiguous
intermediate states remain intermediate after eight attempts (success
$0.3$--$0.6$), a band reproduced in the independent second half of
attempts.
\item \textbf{Early internal signals carry no detectable outcome
information beyond difficulty.} A pre-registered, single-shot test on a
held-out split finds no detectable gain at the frozen forecast point, and
a labeled post-hoc sweep finds none at any window from 128 to 2,048
tokens (Section~\ref{sec:results-prediction}); on public data, a
trace-blind difficulty proxy reaches 0.873 on 192K DeepSeek-R1
generations --- inside the published probe range --- and the closest
published early-window positive replicates at $t{=}4$ (pooled 0.849)
while within problem it is statistically indistinguishable from chance at
all ten anchors (0.496 at $t{=}4$).
\end{itemize}

The claims are scoped deliberately. We study two small open-weight models
(one an explicitly reasoning-tuned release, both run in thinking mode)
plus public corpora from the R1 family; we do not claim internal states
never carry within-attempt information --- at intermediate-answer
positions, in larger models, or in other domains the verdict may differ. What we do claim is that the two controls change
conclusions wherever we have applied them, that they are cheap (our public
data analyses required no generation at all), and that trajectory-level
claims about reasoning should be reported with them. Every
decision rule in this study was frozen in committed amendments before the
outcomes it governs were observed, gate failures included
(Section~\ref{sec:method}, App.~\ref{app:amendments}). Code, the frozen
amendments, and all result artifacts are available at
\url{https://github.com/bulutyigit/problem-not-path}.
